% BHU PYQ -- Manually transcribed & error-corrected
\documentclass[11pt,a4paper]{article}

\usepackage[a4paper,top=2.5cm,bottom=2.5cm,left=2.8cm,right=2.8cm,headheight=14pt]{geometry}
\usepackage{amsmath,amssymb}
\usepackage[T1]{fontenc}
\usepackage[utf8]{inputenc}
\usepackage{lmodern}
\usepackage{microtype}
\usepackage{fancyhdr}
\usepackage{enumitem}
\usepackage{hyperref}
\usepackage{lastpage}
\usepackage{parskip}

\hypersetup{colorlinks=true,urlcolor=black,linkcolor=black,
  pdftitle={CHB-601: Analytical Chemistry-II -- BHU PYQ},pdfauthor={Banaras Hindu University}}

\pagestyle{fancy}\fancyhf{}
\renewcommand{\headrulewidth}{0.4pt}\renewcommand{\footrulewidth}{0.4pt}
\fancyhead[L]{\small   \textit{Banaras Hindu University}}
\fancyhead[R]{\small   \textit{CHB-601: Analytical Chemistry-II}}
\fancyfoot[C]{\small Page  \thepage\ of \pageref{LastPage}}


\newlist{parts}{enumerate}{2}
\setlist[parts,1]{label=(\alph*),leftmargin=*,itemsep=5pt,topsep=3pt}
\setlist[parts,2]{label=(\roman*),leftmargin=*,itemsep=3pt,topsep=2pt}

\newcommand{\pts}[1]{\hfill   \textit{[#1]}}


\begin{document}


\begin{center}
  {\large\bfseries BANARAS HINDU UNIVERSITY}\\[2pt]
  {
\normalsize B.Sc. (Hons.) Semester VI Examination, 2022-23}\\[6pt]
  \rule{0.8\linewidth}{0.5pt}\\[6pt]
  {\large\bfseries Paper No. CHB-601: Analytical Chemistry-II}\\[4pt]
  {
\normalsize Time: 3 Hours\hfill Maximum Marks: 70}
\end{center}
\rule{\linewidth}{0.4pt}
\small



\noindent   \textit{Instructions: Answer five questions in total, including Question~1 which is compulsory.}

\normalsize
\bigskip




\noindent   \textbf{Question 1.} Answer all parts of the following: \pts{5 $ \times$ 2 = 10}

\begin{parts}
  \item List various laboratory neutron sources. How is sensitivity affected by the neutron flux of the source?
  \item Explain the term: reverse-phase chromatography.
  \item Differentiate between cation exchanger and anion exchanger resins.
  \item Write the relationship between radioactive half-life ($t_{1/2}$) and decay constant ($\lambda$).
  \item What is the role of a monochromator in a spectrophotometer?
\end{parts}

\medskip\rule{\linewidth}{0.2pt}




\noindent   \textbf{Question 2.}

\begin{parts}
  \item A solution of organic acid has distribution coefficient $K_D = 10.0$ between ether and water. If $10.0\, \text{g}$ of acid is dissolved in $1.0\, \text{L}$ of water, compare the amounts of acid extracted by: (i) single extraction with $1.0\, \text{L}$ of ether, and (ii) two successive extractions with $500\, \text{mL}$ portions of ether. \pts{8}
  \item Discuss the isotope dilution method applied in radiochemical analysis. Mention its typical applications. \pts{6}
\end{parts}

\medskip\rule{\linewidth}{0.2pt}




\noindent   \textbf{Question 3.}

\begin{parts}
  \item Discuss the principles and applications of Neutron Activation Analysis (NAA). \pts{6}
  \item Calculate the percentage distribution of a substance after three transfers in a Craig counter-current apparatus if the volume of the stationary phase is $4.0\, \text{mL}$, the volume of mobile phase is $8.0\, \text{mL}$, and the distribution coefficient ($K_D$) is $6.0$. \pts{8}
\end{parts}

\medskip\rule{\linewidth}{0.2pt}




\noindent   \textbf{Question 4.}

\begin{parts}
  \item State and derive Beer-Lambert's law. Describe and compare different reasons for deviations from Beer's law (chemical, instrumental, and real deviations). \pts{8}
  \item Potassium dichromate ($ \text{K}_2 \text{Cr}_2 \text{O}_7$) aqueous solution does not obey Beer's law upon dilution. Comment and suggest how you would make this solution obey Beer's law. \pts{6}
\end{parts}

\medskip\rule{\linewidth}{0.2pt}




\noindent   \textbf{Question 5.} Answer any four of the following: \pts{4 $ \times$ 3.5 = 14}

\begin{parts}
  \item Describe the procedure to determine the $ \text{p}K_a$ of an acid-base indicator spectrophotometrically.
  \item Prove that the minimum relative error in concentration in spectrophotometric measurements occurs when $T = 36.8\%$.
  \item Discuss paper chromatography and HPLC separation mechanisms.
  \item Compare adsorption chromatography and partition chromatography.
  \item Draw the block diagram showing the instrumental design of a double-beam spectrophotometer.
\end{parts}

\vfill
\rule{\linewidth}{0.4pt}

\begin{center}\small
  \href{https://bhu-exam.vercel.app/}{Click here to visit our website}\\[4pt]
  \href{https://me-aryan.vercel.app/}{Click here to visit our contributor}\\[4pt]
  \href{https://quashan-me.vercel.app/}{Click here to visit our contributor}
\end{center}

\end{document}
